% Work-MAT | 2. rocnik | hodina 19 | revize 2026-11-09
\documentclass[11pt,a4paper]{article}
\input{../../sablony/pripravy/priprava.sty}
\usepackage{tikz}
\usetikzlibrary{arrows.meta}
\nastavPripravu{2. ročník}{Exponenciální rovnice}{19}{Exponenciální rovnice IV: složitější vytýkání}
\begin{document}
\poznamka{Navazujeme na hodinu 18: nyní bude společný člen nutné nejprve rozpoznat po úpravě více exponentů, někdy se dvěma různými základy. Pořád však řešíme výhradně vytýkáním, nikoli substitucí. Několik typů odpovídá učitelem zaslanému listu „5.1 Exponenciální rovnice“.}
\blok{Vytýkáme na obou stranách rovnice}
U obtížnější rovnice může být součet mocnin na obou stranách. V každé skupině nejprve vytkneme základní mocninu; dostaneme dva jednoduché exponenciální výrazy. Je-li po vydělení jejich poměr opět mocninou stejného základu, můžeme porovnat exponenty. Dělení mocninou $a^x$ je přípustné, protože při $a>0$ nikdy není nulová.
\begin{prikladReseniBox}[title={Příklad 1}]
Vyřeš
\[
3^x+3^{x+1}=7\cdot4^x-4^{x+1}.
\]
\tcblower
\[
\begin{aligned}
3^x(1+3)&=4^x(7-4),\\
4\cdot3^x&=3\cdot4^x,\\
\left(\frac{3}{4}\right)^x&=\frac{3}{4},\\
\left(\frac{3}{4}\right)^x&=\left(\frac{3}{4}\right)^1,\\
x&=1.
\end{aligned}
\]
Zkouška: $3+9=12$ a $7\cdot4-16=12$.
\end{prikladReseniBox}
\begin{prikladReseniBox}[title={Příklad 2}]
Vyřeš rovnici obsahující $3^x$ a $9^x$:
\[
3^x+\frac{9^x}{3}=3^{x+1}+\frac{9^x}{9}.
\]
\tcblower
\[
\begin{aligned}
3^x-3^{x+1}&=\frac{9^x}{9}-\frac{9^x}{3},\\
-2\cdot3^x&=-\frac{2}{9}\,9^x,\\
9\cdot3^x&=9^x,\\
3^{x+2}&=3^{2x},\\
x+2&=2x,\\
x&=2.
\end{aligned}
\]
Dělení výrazem $3^x$ je dovoleno, protože je kladný.
\end{prikladReseniBox}
\begin{prikladReseniBox}[title={Příklad 3}]
Vyřeš $2^{x+3}-3\cdot2^{x+1}=3^{x+1}$.
\tcblower
\[
\begin{aligned}
(8-6)2^x&=3\cdot3^x,\\
2\cdot2^x&=3\cdot3^x,\\
\left(\frac{2}{3}\right)^x&=\frac{3}{2},\\
\left(\frac{2}{3}\right)^x&=\left(\frac{2}{3}\right)^{-1},\\
x&=-1.
\end{aligned}
\]
$K=\{-1\}$.
\end{prikladReseniBox}
\blok{Vytýkání po úpravě exponentů}
Někdy je výhodné přepsat také záporný exponent. Například $4^{-x+3}=4^3\cdot4^{-x}$. Potom vytkneme $4^{-x}$ a zůstane rovnice s jedinou mocninou.
\begin{prikladReseniBox}[title={Příklad 4}]
Vyřeš
\[
7\cdot4^{-x+2}=3\cdot4^{-x+3}-5.
\]
\tcblower
\[
\begin{aligned}
7\cdot16\cdot4^{-x}&=3\cdot64\cdot4^{-x}-5,\\
112\cdot4^{-x}&=192\cdot4^{-x}-5,\\
80\cdot4^{-x}&=5,\\
4^{-x}&=\frac{1}{16}=4^{-2},\\
x&=2.
\end{aligned}
\]
Zkouška: $7\cdot4^0=7$ a $3\cdot4^1-5=7$.
\end{prikladReseniBox}
\begin{prikladReseniBox}[title={Příklad 5}]
Vyřeš $2^{x+3}+2^{x+1}=5\cdot4^{x-1}$.
\tcblower
\[
\begin{aligned}
(8+2)2^x&=\frac{5}{4}\cdot2^{2x},\\
10\cdot2^x&=\frac{5}{4}\cdot(2^x)^2,\\
10&=\frac{5}{4}\cdot2^x,\\
2^x&=8=2^3,\\
x&=3.
\end{aligned}
\]
Vydělili jsme kladným $2^x$, takže nevypadne žádný kořen. Nezavádíme novou proměnnou.
\end{prikladReseniBox}
\begin{prikladReseniBox}[title={Příklad 6}]
Rozhodni, kolik řešení má rovnice $5^{x+2}-25\cdot5^x=c$ v závislosti na konstantě $c$.
\tcblower
\[
5^{x+2}-25\cdot5^x=5^x(25-25)=0.
\]
Pro $c=0$ platí rovnost pro všechna $x\in\mathbb R$. Pro $c\ne0$ nemá rovnice žádné řešení.
\end{prikladReseniBox}
\poznamka{Před přechodem na mocninu podílu vždy zdůvodnit dělení kladnou mocninou. V této hodině se nesnažíme nahradit $2^x$ písmenem $t$; to je náplň dalších dvou hodin.}
\begin{zadaniReseniBox}[title={Samostatná práce 1}]
\begin{enumerate}[label=\alph*)]
\item $5^{x+1}-5^{x-1}=24$.
\item $2^{x+2}+2^{x-2}=17$.
\item $3^x+3^{x+1}=7\cdot4^x-4^{x+1}$.
\item $3^x+\frac{9^x}{3}=3^{x+1}+\frac{9^x}{9}$.
\item $2^{x+3}+2^{x+1}=5\cdot4^{x-1}$.
\item $2^{x+3}-3\cdot2^{x+1}=3^{x+1}$.
\item $7\cdot4^{-x+2}=3\cdot4^{-x+3}-5$.
\item $2^{x+2}+2^{x+1}=4\cdot3^x$.
\item $5^{x+2}-25\cdot5^x=1$.
\item $5^{x+2}-25\cdot5^x=0$.
\end{enumerate}
\tcblower
\begin{enumerate}[label=\alph*)]
\item $\frac{24}{5}\cdot5^x=24$, tedy $x=1$.
\item $\frac{17}{4}\cdot2^x=17$, tedy $x=2$.
\item $4\cdot3^x=3\cdot4^x$, tedy $x=1$.
\item $9\cdot3^x=9^x$, tedy $x=2$.
\item $10\cdot2^x=\frac{5}{4}\cdot2^{2x}$, tedy $x=3$.
\item $2\cdot2^x=3\cdot3^x$, tedy $x=-1$.
\item $80\cdot4^{-x}=5$, tedy $x=2$.
\item $6\cdot2^x=4\cdot3^x$, tedy $\left(\frac{2}{3}\right)^x=\frac{2}{3}$, $x=1$.
\item Levá strana je vždy 0, proto $K=\varnothing$.
\item Jde o identitu: $K=\mathbb R$.
\end{enumerate}
\end{zadaniReseniBox}
\begin{zadaniReseniBox}[title={Rozšíření 1}]
\begin{enumerate}[label=\alph*)]
\item Najdi všechna $x\in\mathbb R$ splňující $2^{x+2}+2^{x+1}=6\cdot3^x$.
\item Vysvětli, proč při řešení neztrácíme kořen dělením.
\end{enumerate}
\tcblower
\begin{enumerate}[label=\alph*)]
\item $6\cdot2^x=6\cdot3^x$, tedy $\left(\frac{2}{3}\right)^x=1=\left(\frac{2}{3}\right)^0$; $x=0$.
\item Obě mocniny jsou kladné pro všechna reálná $x$.
\end{enumerate}
\end{zadaniReseniBox}
\begin{zadaniReseniBox}[title={Domácí úkol}]
\begin{enumerate}[label=\alph*)]
\item $3^x+3^{x+1}=108$.
\item $2^{x+1}+2^{x-1}+2^{x+3}=\frac{21}{8}$.
\item $2^{x+3}-3\cdot2^{x+1}=3^{x+1}$.
\item $5^{x+1}-5^{x-1}=24$.
\item $3^x+3^{x+1}=7\cdot4^x-4^{x+1}$.
\item $3^x+\frac{9^x}{3}=3^{x+1}+\frac{9^x}{9}$.
\item $7\cdot4^{-x+2}=3\cdot4^{-x+3}-5$.
\item $2^{x+3}+2^{x+1}=5\cdot4^{x-1}$.
\item Vysvětli, proč je dělení výrazem $a^x$ přípustné při $a>0$, $a\ne1$.
\item Pro jaké konstanty $c$ má $5^{x+2}-25\cdot5^x=c$ řešení?
\end{enumerate}
\tcblower
\begin{enumerate}[label=\alph*)]
\item $4\cdot3^x=108$, tedy $x=3$.
\item $\frac{21}{2}\cdot2^x=\frac{21}{8}$, tedy $x=-2$.
\item $2\cdot2^x=3\cdot3^x$, tedy $x=-1$.
\item $\frac{24}{5}\cdot5^x=24$, tedy $x=1$.
\item $4\cdot3^x=3\cdot4^x$, tedy $x=1$.
\item $9\cdot3^x=9^x$, tedy $x=2$.
\item $80\cdot4^{-x}=5$, tedy $x=2$.
\item $10\cdot2^x=\frac{5}{4}\cdot2^{2x}$, tedy $x=3$.
\item Platí $a^x>0$ pro všechna reálná $x$; nedělíme nulou.
\item Levá strana je stále 0, tedy pro $c=0$ všechna reálná $x$, pro $c\ne0$ žádné.
\end{enumerate}
\end{zadaniReseniBox}
\poznamka{Domácí úkol asi 25--30 minut, lze vybrat příklady podle rychlosti třídy. Další výukové hodiny 20 a 21 doporučuji věnovat substituci; písemný test T2 je nutné časově posunout až za ně.}
\end{document}
